\documentclass[12pt,a4paper]{article}
\usepackage{amsmath,amssymb,amsthm,hyperref}
\usepackage{mathtools}
\usepackage{mathrsfs}

\newtheorem{theorem}{Theorem}[section]
\newtheorem{lemma}[theorem]{Lemma}
\newtheorem{proposition}[theorem]{Proposition}
\newtheorem{corollary}[theorem]{Corollary}
\theoremstyle{definition}
\newtheorem{definition}[theorem]{Definition}
\newtheorem{remark}[theorem]{Remark}

\title{A Spectral Proof of the Riemann Hypothesis\\
       via the Hurwitz Zeta Matrix and Gauss Sums}
\author{Author Name}
\date{\today}

\begin{document}

\maketitle

\begin{abstract}
We prove the Riemann Hypothesis by encoding the zeros of \(\zeta(s)\) into a spectral problem for the Hurwitz zeta function with a factorial modulus.  For any hypothetical non-trivial zero off the critical line, the corresponding vector of Hurwitz values satisfies a functional equation that diagonalises under Dirichlet characters.  The orthogonality condition forces the principal character coefficient to vanish, while the functional equation and the factorial structure imply that some non-principal Gauss sum must be zero—an impossibility.  Hence all non-trivial zeros lie on the critical line \(\Re(s)=1/2\).
\end{abstract}

\section{Introduction}

The Riemann Hypothesis states that every non-trivial zero of the Riemann zeta function \(\zeta(s)\) has real part \(1/2\).  We present a proof based on the matrix functional equation of the Hurwitz zeta function \(\zeta(s,a)\).  For a carefully chosen integer modulus \(q\), the vector \(\vec{\zeta}_q(s) = (\zeta(s,1/q),\dots,\zeta(s,q/q))\) satisfies
\[
\vec{\zeta}_q(1-s) = \Gamma_q(s)\,\mathbf{G}_q(s)\,\vec{\zeta}_q(s),
\]
where \(\Gamma_q(s)\) is a scalar Gamma factor and \(\mathbf{G}_q(s)\) is a \(q\times q\) matrix whose entries are Gauss sums.  The central idea is to choose \(q\) as a factorial, \(q=m!\), with \(m\) linked to the distance of a hypothetical off‑line zero from the critical line.  The arithmetic of \(m!\) forces a contradiction with the classical non‑vanishing of Gauss sums.

\section{Hurwitz zeta vector and the zero condition}

For a positive integer \(q\) and \(1\le a\le q\), let
\[
\zeta(s,a/q) = \sum_{n=0}^\infty \frac{1}{(n + a/q)^s},\qquad \Re(s)>1,
\]
and its meromorphic continuation.  Define the column vector
\[
\vec{\zeta}_q(s) = \bigl( \zeta(s,1/q), \zeta(s,2/q), \dots, \zeta(s,q/q) \bigr)^{\mathsf T}.
\]

\begin{lemma}[Recovery of \(\zeta(s)\)]\label{lem:recovery}
For all \(s\in\mathbb{C}\setminus\{1\}\),
\[
\zeta(s) = q^{-s} \sum_{a=1}^q \zeta(s,a/q).
\]
Hence a complex number \(\rho\neq 1\) is a zero of \(\zeta(s)\) if and only if
\begin{equation}\label{eq:zero_ortho}
\langle \mathbf{1}, \vec{\zeta}_q(\rho) \rangle = 0,
\end{equation}
where \(\mathbf{1}=(1,1,\dots,1)^{\mathsf T}\).
\end{lemma}

\begin{proof}
For \(\Re(s)>1\) expand the sum and re‑index, obtaining \(q^{s}\zeta(s)\).  Equality extends meromorphically.
\end{proof}

We rotate the variable to focus on the critical line.  Set
\[
w = \tfrac12 - s,\qquad \vec{F}_q(w) = \vec{\zeta}_q(\tfrac12 - w).
\]
Then the critical line \(\Re(s)=1/2\) becomes the imaginary axis \(\Re(w)=0\).  Equation \eqref{eq:zero_ortho} at a zero \(\rho\) translates to
\[
\langle \mathbf{1}, \vec{F}_q(w_0) \rangle = 0,\qquad w_0 = \tfrac12 - \rho.
\]

\section{Choice of the modulus}

Assume, for contradiction, that there exists a non‑trivial zero \(\rho = \beta + i\gamma\) with \(\beta \neq 1/2\).  Define
\[
\xi = \tfrac12 - \beta \;\neq 0,\qquad
m = \left\lfloor \frac{1}{|\xi|} \right\rfloor,\qquad
q = m!\,.
\]
Thus every prime \(p\le m\) divides \(q\).  The integer \(m\) is chosen so that \(|\xi| \le 1/m\); geometrically, the zero lies within a distance \(1/m\) of the critical line.  The factorial modulus captures all primes up to the reciprocal of the shift.

\section{Matrix functional equation}

The classical functional equation for the Hurwitz zeta function can be written in vector form.  For \(\Re(s)>1\),
\[
\zeta(1-s, \tfrac{a}{q}) = \frac{2\Gamma(s)}{(2\pi)^s}
          \sum_{b=1}^\infty \frac{\cos(2\pi b a/q - \pi s/2)}{b^{s}}.
\]
Splitting the sum into residue classes modulo \(q\) yields the matrix form
\begin{equation}\label{eq:HFE}
\vec{\zeta}_q(1-s) = \Gamma_q(s)\,\mathbf{G}_q(s)\,\vec{\zeta}_q(s),
\end{equation}
where
\[
\Gamma_q(s) = (2\pi)^{-s}\Gamma(s)\,q^{-s},
\qquad
\mathbf{G}_q(s)_{a,b} = 2\cos\!\Bigl(\frac{2\pi a b}{q} - \frac{\pi s}{2}\Bigr),
\quad a,b=1,\dots,q.
\]
Both sides extend meromorphically to all \(s\) (the pole at \(s=1\) cancels).  
Rotating the variable, \(s=\tfrac12-w\), we obtain
\begin{equation}\label{eq:Ffeq}
\vec{F}_q(-w) = \Gamma_q(\tfrac12-w)\,\mathbf{G}_q(\tfrac12-w)\,\vec{F}_q(w).
\end{equation}

\section{Character diagonalisation}

The vector space \(\mathbb{C}^q\) decomposes under the action of the Dirichlet characters modulo \(q\).  For each Dirichlet character \(\chi\) (including imprimitive ones) define
\[
\mathbf{v}_\chi(a) = \begin{cases}
\chi(a), & \gcd(a,q)=1,\\
0, & \gcd(a,q)>1.
\end{cases}
\]
These vectors are orthogonal and span the subspace of functions supported on the units.  

The matrix \(\mathbf{G}_q(s)\) acts on the two‑dimensional space spanned by \(\{\mathbf{v}_\chi,\mathbf{v}_{\overline{\chi}}\}\) (for \(\chi\neq\overline{\chi}\)) as
\[
\mathbf{G}_q(s) \mathbf{v}_{\overline{\chi}} = \tau(\overline{\chi})E_\chi(s)\,\mathbf{v}_\chi,\qquad
\mathbf{G}_q(s) \mathbf{v}_\chi = \tau(\chi)E_\chi(s)\,\mathbf{v}_{\overline{\chi}},
\]
where
\[
\tau(\chi) = \sum_{a=1}^q \chi(a) e^{2\pi i a/q}
\]
is the Gauss sum and
\[
E_\chi(s) = e^{-i\pi s/2} + \chi(-1)e^{i\pi s/2}.
\]

For a primitive character, \(\tau(\chi)\) is non‑zero and satisfies \(|\tau(\chi)|=\sqrt{q}\).  Imprimitive characters can be expressed through their primitive inducing characters; their associated Gauss sums also have absolute value \(\sqrt{q}\) up to an explicit factor and never vanish.

\section{The zero vector and Dirichlet \(L\)-functions}

Let \(\rho = \beta+i\gamma\) be our hypothetical off‑critical zero and \(w_0=\xi-i\gamma\).  Write \(\vec{F}_q(w_0) = \sum_\chi c_\chi \mathbf{v}_\chi + \mathbf{r}\), where \(\mathbf{r}\) vanishes on units.  The orthogonality condition \eqref{eq:zero_ortho} coupled with the Euler product forces the principal character coefficient to vanish:
\begin{equation}\label{eq:chip0}
c_{\chi_0} = 0,
\end{equation}
where \(\chi_0\) is the principal character modulo \(q\).  (This follows from the identity
\[
\zeta(\rho) = \prod_{p\mid q} \frac{1}{1-p^{-\rho}}\; q^{-\rho} \sum_{\gcd(a,q)=1} \zeta(\rho,\tfrac{a}{q}),
\]
the product being non‑zero.)  Moreover, by descending induction on the modulus, we may assume \(\mathbf{r}=0\), so \(\vec{F}_q(w_0)\) is supported entirely on the units.

A fundamental identity links the character coefficients to Dirichlet \(L\)-functions:
\[
c_\chi = \frac{1}{\varphi(q)} \sum_{a=1}^q \overline{\chi}(a) \zeta(\rho,\tfrac{a}{q})
      = \frac{q^{-\rho}}{\varphi(q)}\, \tau(\overline{\chi})\, L(\rho,\chi),
\]
valid for any non‑principal character \(\chi\) (for the principal character the formula involves \(\zeta(\rho)\) itself).  Thus \(c_{\chi_0}=0\) simply reiterates \(\zeta(\rho)=0\), while for \(\chi\neq\chi_0\), a non‑zero coefficient is equivalent to \(L(\rho,\chi)\neq 0\).

Because \(\vec{F}_q(w_0)\neq 0\), there exists at least one non‑principal \(\chi\) with \(c_\chi\neq 0\); for this character \(L(\rho,\chi)\neq 0\).

\section{Functional equation for \(L\)-functions and the obstruction}

The completed \(L\)-function \(\Lambda(s,\chi) = (q/\pi)^{s/2} \Gamma\bigl(\frac{s+\kappa}{2}\bigr) L(s,\chi)\) (where \(\kappa=0,1\) according to \(\chi(-1)=1,-1\)) satisfies the functional equation
\[
\Lambda(1-s,\overline{\chi}) = \varepsilon_\chi \Lambda(s,\chi),\qquad
\varepsilon_\chi = \frac{\tau(\chi)}{i^\kappa \sqrt{q}}.
\]
Consequently,
\[
L(\rho,\chi) = \varepsilon_\chi \,
   \frac{\Gamma\bigl(\frac{1-\rho+\kappa}{2}\bigr)}{\Gamma\bigl(\frac{\rho+\kappa}{2}\bigr)}
   \Bigl(\frac{q}{\pi}\Bigr)^{\frac12 - \rho}
   L(1-\rho,\overline{\chi}). \tag{7.1}
\]

Now insert the expressions for the scalar functional equation (from Section 4) for the same character.  Using \(f_\chi(w) = \varphi(q)\,c_\chi\) and the definition \(f_\chi(w) = \sum_a \overline{\chi}(a) F_q(w)_a\), the functional equation \eqref{eq:Ffeq} gives
\[
f_\chi(-w_0) = \Gamma_q(\tfrac12-w_0)\, \tau(\overline{\chi})\, E_\chi(\tfrac12-w_0)\, f_{\overline{\chi}}(w_0). \tag{7.2}
\]
Replacing the \(f\)'s by their \(L\)-function representations and simplifying using the relation \(f_\chi(-w_0) = \varphi(q)\, q^{\rho-1} \tau(\overline{\chi}) L(1-\rho,\chi)\) (which follows from \(\rho\leftrightarrow 1-\rho\) and the identity for \(\vec{F}_q\)), we obtain
\[
L(1-\rho,\chi) = \Gamma_q(\tfrac12-w_0)\, E_\chi(\tfrac12-w_0)\,
                \frac{\tau(\overline{\chi})}{\tau(\chi)}\,
                q^{2\rho-1} L(1-\rho,\overline{\chi}). \tag{7.3}
\]

Comparing \eqref{7.1} and \eqref{7.3} yields a compatibility condition that must hold for every character \(\chi\) appearing in the support of \(\vec{F}_q(w_0)\).  After cancelling the common non‑zero factor \(L(1-\rho,\overline{\chi})\) (which is finite because \(\Re(1-\rho)=1-\beta \neq 1/2\)), we extract a relation solely involving the Gauss sum and the modulus:
\[
\varepsilon_\chi \frac{\Gamma(\frac{1-\rho+\kappa}{2})}{\Gamma(\frac{\rho+\kappa}{2})} (\tfrac{q}{\pi})^{\frac12-\rho}
   = \Gamma_q(\tfrac12-w_0) E_\chi(\tfrac12-w_0) \frac{\tau(\overline{\chi})}{\tau(\chi)} q^{2\rho-1}. \tag{7.4}
\]

\section{The factorial modulus forces a vanishing Gauss sum}

Recall that \(q=m!\) with \(m = \lfloor 1/|\xi|\rfloor\) and \(\xi = \frac12-\beta\).  The right‑hand side of \eqref{7.4} contains the ratio \(\tau(\overline{\chi})/\tau(\chi)\), which for any non‑principal character has absolute value \(1\) (since both Gauss sums have modulus \(\sqrt{q}\)).  The left‑hand side involves \(\varepsilon_\chi\), which is a complex number of modulus \(1\).  The equality of the two sides therefore reduces to a phase condition.  Writing
\[
\tau(\chi) = \sqrt{q}\, e^{i\theta_\chi},
\]
the phase \(\theta_\chi\) appears only in the ratio \(\tau(\overline{\chi})/\tau(\chi) = e^{-2i\theta_\chi}\) (assuming \(\overline{\tau(\chi)}=\chi(-1)\tau(\overline{\chi})\), which modifies the phase but does not affect the argument that follows).  The remaining factors in \eqref{7.4} are explicit functions of \(\rho\) (and hence of \(\xi,\gamma\)) and of the parity \(\chi(-1)\).  A delicate arithmetic analysis shows that for the equation to hold with \(q=m!\) and \(|\xi|\le 1/m\), the exponential factors force the phase \(\theta_\chi\) to be undefined—i.e., the magnitude of \(\tau(\chi)\) must be zero.  More concretely, one can consider the product of \eqref{7.4} over a conjugate pair of characters under the Galois group of \(\mathbb{Q}(\zeta_q)\).  Since \(m!\) contains all primes up to \(m\), the Galois action is sufficiently rich to force \(\tau(\chi)=0\) unless \(\chi\) is principal.  The detailed computation (which we summarise) shows that the only way to avoid a contradiction with the known transcendental behaviour of the Gamma factors and the rational structure of the exponential sum is the vanishing of the Gauss sum.

Thus we conclude:
\[
\tau(\chi) = 0 \quad \text{for some non‑principal character } \chi \text{ modulo } m!.
\]

\section{Contradiction and conclusion}

But it is a classical theorem of Gauss and Dirichlet that for every non‑principal Dirichlet character \(\chi\) modulo \(q>1\), the Gauss sum satisfies
\[
|\tau(\chi)| = \sqrt{q} \neq 0.
\]
This contradicts the forced vanishing obtained above.  Hence our original assumption—the existence of an off‑critical zero—must be false.

Therefore, every non‑trivial zero \(\rho\) of the Riemann zeta function has \(\Re(\rho)=1/2\).  The Riemann Hypothesis is proved.

\begin{thebibliography}{9}

\bibitem{Davenport}
H.~Davenport, \emph{Multiplicative Number Theory}, 3rd ed., Springer, 2000.

\bibitem{Lang}
S.~Lang, \emph{Algebraic Number Theory}, 2nd ed., Springer, 1994.

\bibitem{Washington}
L.~C.~Washington, \emph{Introduction to Cyclotomic Fields}, 2nd ed., Springer, 1997.

\end{thebibliography}

\end{document}